\documentclass[HARVARD,Utopia2COL]{WileyNJDv5}

\usepackage{amsmath,amssymb,booktabs}
\usepackage{graphicx}
\hypersetup{hidelinks}

\articletype{Research Article}
\journal{Journal of Futures Markets}
\raggedbottom

\begin{document}

\title{When Does Posterior Integration Change Option Values? Evidence from WTI Average Price Options}

\author[1]{Heriberto Espino Montelongo}
\authormark{ESPINO MONTELONGO}
\titlemark{POSTERIOR INTEGRATION IN WTI AVERAGE PRICE OPTIONS}

\address[1]{\orgname{Universidad de las Am\'ericas Puebla}, \orgaddress{\state{Puebla}, \country{Mexico}}}
\corres{Corresponding author: Heriberto Espino Montelongo. \email{[corresponding email]}}

\abstract[Abstract]{
This paper studies a narrow Bayesian option-pricing question: when does integrating a posterior for historically estimated volatility change a WTI Average Price Option point value relative to plugging in the posterior mean, and how large is that correction relative to changing the volatility input itself? Statistical inference is conducted under the physical measure and valuation under the risk-neutral measure. A 175,000-history synthetic mechanism map shows that the posterior-integrated minus posterior-mean correction is governed by posterior dispersion and volatility curvature; the largest observed absolute gap is 0.1384, and the second-order term $\tfrac12 C^{Q\prime\prime}(\bar\sigma)\operatorname{Var}(\sigma\mid\mathcal D)$ closely tracks the realized gap. In the October-2026 WTI panel, posterior-integrated and posterior-mean RMSE are 0.5177 and 0.5183. Their mean absolute price difference is 0.000745, while the mean width of the 95\% posterior model-price interval is 0.2733, showing that a small mean-price correction does not imply negligible parameter-induced price uncertainty. In a strict forward-in-time panel spanning seven expiries and 2,381 holdouts, the long-window historical-volatility baseline has RMSE 3.4090, compared with 0.1725 for an expanding prior-date APO smile and 0.1594 for a previous-day smile. An independent vanilla-WTI validation gives RMSE 0.5562 for the historical baseline, 0.2599 for a previous-day LO surface, and 0.2686 for the corresponding APO smile on 253 common-support observations. The empirical result is therefore interpreted as an option-informed versus long-window historical-volatility comparison, not as a pure structural decomposition between $\mathbb P$ and $\mathbb Q$. Posterior integration can nevertheless become appreciable in weak-information, high-curvature states.
}

\keywords{Bayesian option pricing; Asian options; WTI crude oil; parameter uncertainty; commodity derivatives}

\maketitle


\section{Introduction}\label{sec:introduction}

Derivative values are conditional on parameters and state variables that are not observed without error. In the Black--Scholes--Merton benchmark, volatility is often inserted into the pricing map as though it were known, even when it has been estimated from a finite return history. A Bayesian treatment replaces that point estimate with a posterior distribution and can propagate posterior uncertainty through the derivative-pricing map. This idea is not new to option pricing: Bayesian learning has been used to explain implied-volatility dynamics and improve out-of-sample option forecasts \cite{guidolin2003}, risk-neutral predictive densities have been constructed while integrating parameter uncertainty \cite{romboutsstentoft2014}, and Bayesian calibration has been used to propagate uncertainty in calibrated model parameters to prices of exotic derivatives \cite{guptareisinger2014}. The question in this paper is therefore narrower: \emph{when does posterior integration materially change a point price, and how important is that effect relative to changing the volatility or pricing specification itself?} The paper does not attempt to quantify every consequence of parameter uncertainty: hedging error, latent-state uncertainty, calibration uncertainty in richer models, and decision-making under posterior loss distributions are outside its scope.

Average-price commodity options provide a useful setting for that comparison. Their payoff depends on a path of underlying prices, so the curvature of the pricing map can change materially with moneyness, maturity, and the fraction of the averaging window already fixed. The numerical literature has long exploited geometric-average structure to price arithmetic-average claims efficiently \cite{kemnavorst1990,curran1994}. Commodity-specific models add mean reversion, convenience yield, stochastic interest rates, and jumps \cite{schwartz1997,ewaldwu2023}. Most importantly for the present application, \cite{shirayatakahashi2011} study WTI average-price options under Heston and extended $\lambda$-SABR dynamics, explicitly calibrating richer risk-neutral models to more liquid American vanilla options on WTI futures before valuing the average option. \cite{ganwangyang2020} provide a distinct model-free machine-learning application using real WTI Average Price Option data. These contributions mean that neither WTI APO pricing nor the use of richer calibrated risk-neutral dynamics is itself the novelty here.

The empirical distinction between parameter uncertainty and model specification is also central in the broader option-pricing literature. Richer stochastic-volatility and jump models can reduce pricing errors, but additional complexity does not improve every performance criterion \cite{bakshicaochen1997,cumminsesposito2025}. Flexible implied-volatility functions can fit the contemporaneous cross-section extremely well and still deteriorate out of sample, which makes forward validation essential \cite{dumasflemingwhaley1998}. In crude oil specifically, option-implied volatility measures contain forward-looking information beyond historical return models \cite{gildertsiaras2020}, commodity variance risk premia imply an economically meaningful distinction between physical and risk-neutral variation \cite{teeting2017}, and volatility-of-volatility risk is itself priced in oil-option returns \cite{roh2021}. These results motivate treating an option-implied scalar volatility as an \emph{effective} or model-equivalent risk-neutral state rather than as a structural diffusion coefficient.

This paper studies that comparison using CME WTI Average Price Options (APOs). The contract pays against the arithmetic average of daily first-nearby WTI futures settlements over a calendar month. Consequently, a contract-consistent valuation exercise must account for the first-nearby roll convention and partial fixing: once the averaging month has begun, some components of the final average are known while the remaining components are stochastic. Treating the option as an ordinary European claim on a single fixed-maturity futures contract would change the object being valued.

The central methodological principle is a strict separation between inference under the physical measure and valuation under the risk-neutral measure. Historical returns are used to infer a data-generating process under $\mathbb P$. Posterior uncertainty in volatility is then mapped into option values under $\mathbb Q$. The physical drift is not inserted into the risk-neutral pricing dynamics. This distinction follows the no-arbitrage separation between physical inference and risk-neutral valuation \cite{blackscholes1973,merton1973}. Under the maintained constant-volatility benchmark, the same diffusion coefficient can coherently be used under both measures; whether that historical volatility is empirically adequate for the option market is a separate specification question.

Let $C^{Q}(\sigma)$ denote the risk-neutral value of an arithmetic-average option conditional on volatility $\sigma$. The posterior-integrated rule is
\begin{equation}
 C_{PI}=\mathbb E\!\left[C^{Q}(\sigma)\mid\mathcal D\right],
 \label{eq:intro_fb}
\end{equation}
whereas the posterior-mean plug-in rule is
\begin{equation}
 C_{PM}=C^{Q}\!\left(\mathbb E[\sigma\mid\mathcal D]\right).
 \label{eq:intro_pm}
\end{equation}
The two need not coincide because the option-pricing map is nonlinear. A second-order expansion around $\bar\sigma=\mathbb E[\sigma\mid\mathcal D]$ gives
\begin{equation}
 C_{PI}-C_{PM}
 \approx
 \frac{1}{2}C^{Q\prime\prime}(\bar\sigma)
 \operatorname{Var}(\sigma\mid\mathcal D),
 \label{eq:intro_taylor}
\end{equation}
which isolates two forces: posterior dispersion and curvature of the pricing function. Parameter uncertainty should matter most where the posterior is diffuse and option value is highly curved in volatility.

The paper develops four layers of evidence. First, repeated-sampling experiments and a 175,000-history mechanism map provide controlled environments in which the true volatility is known. The mechanism map varies historical sample size from 21 to 1,260 observations, true volatility from 0.10 to 0.80, moneyness from 0.60 to 1.50, maturity from 21 to 252 days, and the fraction of the averaging window already fixed. This makes it possible to identify where the approximation in Equation~\eqref{eq:intro_taylor} succeeds and where posterior integration can become economically material.

Second, the paper develops a contract-consistent empirical application using historical WTI APO settlements obtained through Barchart. The downloaded database contains 213 unique call and put contracts and 11,179 option-date observations across expiries from September 2026 through June 2029. The October-2026 baseline contains 310 retained option-date observations across 12 valuation dates from August 25 through September 10, 2026. Additional September and October panels provide 402 contract-date observations for implied-volatility and heavy-tail robustness analyses. In the Barchart histories used here, \texttt{Latest} is the CME settlement field; the source-field name is retained for provenance and is not interpreted as an intraday last trade.

Third, the paper distinguishes posterior integration from volatility-input choice using the APO market itself. It inverts the maintained pricing map to obtain Curran-equivalent APO-implied volatilities and then asks whether volatility smiles estimated only from earlier dates improve subsequent APO valuation. This strict forward-in-time design spans seven expiries and prevents same-day target prices from entering the fitted smile. Pre-specified rolling and exponentially weighted historical-volatility comparators test whether recency alone explains the difference. A horizon-matched Gaussian GARCH(1,1) forecast then tests whether conditional heteroskedasticity and maturity-specific historical variance forecasting materially narrow the gap, while a same-contract prior-IV carry benchmark tests whether flexible cross-sectional smile fitting is necessary.

Fourth, the paper breaks that remaining within-family link using standard monthly WTI futures options. Official prior-date LO option settlements and matching CL futures settlements are inverted through an American futures-option model to obtain an independent risk-neutral volatility surface. That surface is mapped to the APO fixing horizon and used to price October-2026 APOs without using APO prices to construct the external state. This cross-instrument experiment asks whether the pricing gain attributed to risk-neutral volatility information survives when the source and target option families are separated.

The main results separate the mean-price correction from both numerical error and posterior price uncertainty. In the October-2026 baseline, posterior-integrated and posterior-mean pricing are almost identical: pooled RMSE is 0.5177 versus 0.5183 across 310 observations. The mean absolute PI--PM difference is 0.000745, while the mean width of the 95\% posterior model-price interval is 0.2733. Thus PI and PM can agree closely even when parameter uncertainty induces a visibly dispersed distribution of theoretical prices. High-precision pseudo-random Monte Carlo, randomized Sobol, and Curran conditioning resolve the PI--PM gap to a scale far below its observed magnitude.

The synthetic mechanism map nevertheless finds states in which posterior integration is appreciable. The largest observed absolute gap is 0.1384, concentrated in short historical samples, high volatility, long remaining horizons, extreme moneyness, and limited realized fixing. The second-order curvature approximation tracks these gaps closely, so the mechanism is not specific to one numerical engine.

The larger empirical difference is between volatility inputs. Across 2,381 strict forward-in-time holdouts spanning seven expiries, the long-window historical-volatility posterior-integrated baseline has MAE 2.6187 and RMSE 3.4090, compared with MAE 0.1088 and RMSE 0.1725 for an expanding prior-date APO smile and MAE 0.1075 and RMSE 0.1594 for a previous-day smile. Recent historical volatility does not close that gap: the best pooled pre-specified recent-history rule, based on the latest 63 returns, has MAE 4.5201 and RMSE 5.6465. A horizon-matched Gaussian GARCH(1,1) forecast also remains near the historical baseline, with MAE 2.6730 and RMSE 3.4441. Conversely, on 2,366 exact common rows, simply carrying each contract's latest strictly prior implied volatility gives MAE 0.0992 and RMSE 0.1403, slightly better than the fitted previous-day smile. Reconstructing the physical return history contract by contract from first-nearby CL settlements changes the long-window historical RMSE only to 3.4451, so the result is not a Yahoo roll artifact.

The independent vanilla-WTI experiment is the strongest check on within-family circularity. On 253 October-2026 contract-dates that lie inside the observed prior LO moneyness support, historical-volatility posterior integration has MAE 0.4464 and RMSE 0.5562, compared with 0.1559 and 0.2599 for the external previous-day LO surface. The corresponding prior-date APO smile has MAE 0.1676 and RMSE 0.2686. Valuation-date cluster bootstraps show a large external-LO improvement over the historical benchmark, while differences between the external LO and prior-date APO surfaces are small.

These comparisons do not identify a pure structural difference between physical and risk-neutral volatility. The recent-history experiment shows that recency by itself does not reproduce the option-informed result, the horizon-matched GARCH experiment shows that a standard dynamic historical forecast does not materially narrow the pooled gap, and the contract-IV carry experiment shows that a flexible fitted smile is not required. The defensible empirical statement is therefore narrower: in this sample, prior option-implied volatility states predict APO settlements much better than the maintained historical-volatility benchmarks, while integrating the historical posterior changes point prices very little.

The paper's contribution is diagnostic rather than a new pricing method. It quantifies the posterior-integration correction, maps the states in which that correction grows, separates it from posterior price dispersion and numerical error, and places it beside an empirically larger volatility-input margin. This provides a practical benchmark for deciding when Bayesian propagation of an estimated scalar volatility is worth carrying before adding further model complexity.

The remainder of the paper is organized as follows. Section~\ref{sec:contract} describes the WTI Average Price Option and the empirical panel. Section~\ref{sec:methodology} presents physical-measure inference, risk-neutral valuation, and the competing pricing rules. Section~\ref{sec:design} describes the synthetic, empirical, within-APO forward, and independent vanilla-WTI validation designs. Section~\ref{sec:results} reports the mechanism and market evidence. Section~\ref{sec:robustness} evaluates numerical accuracy, posterior calibration, heavy-tailed returns, liquidity dependence, and the interpretation of option-informed risk-neutral volatility. Section~\ref{sec:conclusion} concludes.

\section{WTI Average Price Options and the empirical setting}\label{sec:contract}

\subsection{Contract mechanics}

CME WTI Average Price Options are cash-settled options whose payoff is linked to the arithmetic average of daily settlements of first-nearby WTI crude-oil futures during the relevant calendar month. For an averaging month with fixing dates $t_1,\ldots,t_M$, define
\begin{equation}
 A_T=\frac{1}{M}\sum_{j=1}^{M}F^{(1)}_{t_j},
 \label{eq:apo_average}
\end{equation}
where $F^{(1)}_{t_j}$ denotes the settlement of the first-nearby CL futures contract applicable to fixing date $t_j$. A call and put have terminal payoffs
\begin{align}
 H_C&=(A_T-K)^+,\\
 H_P&=(K-A_T)^+.
\end{align}
The first-nearby designation matters because the futures contract contributing to the average can change during the averaging period. The payoff is therefore not, in general, an average of repeated observations of one fixed futures maturity.

At an option valuation date $t$ inside the averaging month, let $\mathcal R_t$ denote fixing dates already realized and $\mathcal U_t$ the remaining fixing dates. The final average can be decomposed as
\begin{equation}
 A_T=\frac{1}{M}
 \left(
 \sum_{j\in\mathcal R_t}F^{(1)}_{t_j}
 +
 \sum_{j\in\mathcal U_t}F^{(1)}_{t_j}
 \right).
 \label{eq:partial_fixing}
\end{equation}
This decomposition is central to the empirical pricing exercise. Realized fixings are known constants; only the remaining components are simulated under the risk-neutral measure. As the averaging month progresses, the option mechanically loses exposure to future volatility because a larger share of $A_T$ has already been fixed.

A useful state variable for cross-sectional analysis is the risk-neutral expected final average,
\begin{equation}
 \widehat A^{Q}_{t,T}=
 \frac{1}{M}
 \left(
 \sum_{j\in\mathcal R_t}F^{(1)}_{t_j}
 +
 \sum_{j\in\mathcal U_t}\mathbb E_t^Q[F^{(1)}_{t_j}]
 \right).
 \label{eq:expected_average}
\end{equation}
We define log moneyness observation by observation as
\begin{equation}
 m_{t,K,T}=\log\!\left(\frac{K}{\widehat A^{Q}_{t,T}}\right).
 \label{eq:moneyness}
\end{equation}
For calls, negative values correspond to strikes below the expected average; for puts the economic interpretation is reversed in terms of intrinsic direction. This measure is preferable to assigning a contract a fixed moneyness category using the futures price on the final sample date.

\subsection{Why WTI is informative for parameter uncertainty}

WTI provides a demanding empirical environment for the question studied in this paper. Oil prices can exhibit rapid changes in level, realized volatility, and the shape of the futures curve. These changes imply that the same option can move through substantially different moneyness states during its observed history. They also create episodes in which a posterior estimated from a shorter historical window is more dispersed than one estimated from a long sample.

The analysis does not attribute individual price moves mechanically to geopolitical events. Such causal statements would require a separate event-study design. Instead, the paper conditions its claims on observed return histories, posterior dispersion, option state variables, and the market-implied volatility information available before each valuation date.

\subsection{Barchart option histories}

The empirical database was assembled from individual Barchart historical series for WTI Average Price Options. The downloaded maturities are September 2026, October 2026, November 2026, March 2027, September 2027, March 2028, September 2028, and June 2029. Calls and puts were retained as available. Rather than choosing only contracts that appear near the money on one date, the raw panel keeps all downloaded contracts and lets moneyness vary by date according to Equation~\eqref{eq:moneyness}.

The ingestion pipeline recursively audits the raw CSV directory, parses option symbols, derives strike and option type, checks the symbol-implied expiry against the storage folder, and de-duplicates contracts before constructing the option-date panel. The current audit identifies 214 input files, 213 unique contracts, and 11,179 option-date observations. One file is stored under a folder inconsistent with its symbol-implied expiry; the parser records the mismatch and uses the contract symbol rather than silently trusting the directory name.

The histories contain end-of-day fields including open, high, low, latest, volume, and open interest. In the source histories used for this study, Barchart \texttt{Latest} is the end-of-day settlement field. At least one contract/date comparison also reproduces the corresponding CME daily settlement exactly. The field is therefore used as the market settlement benchmark, while the paper stops short of claiming that the full proprietary panel has been independently reconciled observation by observation against official CME bulletins. This distinction matters because many histories contain zero reported volume while open interest remains positive.

Liquidity is consequently used as a diagnostic and robustness dimension rather than a binary requirement that every observation have positive trading volume. Results are reported for the full eligible panel, positive-volume subsets, and multiple open-interest thresholds. Contracts with sparse trading can still contain an exchange settlement, but conclusions that depend exclusively on such observations are treated cautiously.

The data pipeline also produces a variable called \texttt{richest}, inherited from an earlier data-audit routine. It measures information completeness using the number of sessions, calendar span, and field completeness. It does \emph{not} measure liquidity, trading intensity, or economic importance. The variable is retained for data-quality diagnostics but is not used as the economic contract-selection rule.

\subsection{Maturity coverage}

The chosen expiries deliberately span the maturity dimension rather than concentrating exclusively on consecutive nearby months. The first three maturities provide short-horizon contracts, while March 2027 through June 2029 extend the cross section toward medium and long horizons. This structure is useful because Equation~\eqref{eq:intro_taylor} predicts that the importance of posterior uncertainty depends not only on posterior dispersion but also on the curvature of the pricing map, which can change materially with time to maturity.

The number of available strikes and the depth of trading are not constant across maturities. In particular, the far end of the curve contains fewer contracts. The empirical design does not force a balanced panel. Instead, it preserves the observed contract set and treats maturity coverage, liquidity, and availability of matching CL futures histories as explicit scope constraints.

\section{Methodology}\label{sec:methodology}

\subsection{Inference under the physical measure}

The baseline statistical model assumes that the observed underlying follows a geometric Brownian motion under the physical measure,
\begin{equation}
 dS_t=\mu S_t\,dt+\sigma S_t\,dW_t^{P}.
 \label{eq:p_gbm}
\end{equation}
For equally spaced observations with interval $\Delta t$, log returns satisfy
\begin{equation}
 R_i=\log\!\left(\frac{S_{t_i}}{S_{t_{i-1}}}\right)
 \sim
 \mathcal N\!\left[
 \left(\mu-\frac{1}{2}\sigma^2\right)\Delta t,
 \sigma^2\Delta t
 \right].
 \label{eq:return_likelihood}
\end{equation}
The likelihood is the product of the conditional return densities. The diffusion is intentionally simple: the paper's primary object is the incremental effect of parameter uncertainty and its interaction with the pricing map, not a claim that constant-volatility GBM is a complete model of crude-oil returns.

The baseline prior specification is
\begin{align}
 \mu&\sim\mathcal N(0,1),\\
 \sigma&\sim\operatorname{InvGamma}(2,0.1),\qquad \sigma>0.
 \label{eq:priors}
\end{align}
Sampling is performed in $(\mu,\eta)$ with $\eta=\log\sigma$ so that positivity is automatic. The transformed log posterior is
\begin{equation}
 \log \pi(\mu,\eta\mid\mathcal D)
 =\log L(\mu,e^\eta;\mathcal D)
 +\log p(\mu)+\log p(e^\eta)+\eta+C,
 \label{eq:posterior}
\end{equation}
where the final $\eta$ is the Jacobian term.

A symmetric random-walk Metropolis proposal is used,
\begin{equation}
 \theta'=\theta+\varepsilon,
 \qquad
 \varepsilon\sim\mathcal N(0,\Sigma_q),
\end{equation}
with $\theta=(\mu,\eta)$. The acceptance probability is
\begin{equation}
 \alpha(\theta,\theta')=
 \min\left\{1,
 \frac{\pi(\theta'\mid\mathcal D)}{\pi(\theta\mid\mathcal D)}
 \right\},
\end{equation}
following the Metropolis--Hastings construction \cite{metropolis1953,hastings1970}. The empirical implementation uses multiple chains and Gelman--Rubin $\hat R$ diagnostics for $\mu$ and $\sigma$. For the large synthetic mechanism and calibration experiments, $\mu$ is also integrated analytically under its Gaussian prior, leaving a one-dimensional posterior in $\sigma$ that is evaluated by dense quadrature. This provides an MCMC-independent numerical reference without changing the statistical model.

\subsection{Risk-neutral valuation}

Historical inference and derivative valuation are kept conceptually and computationally separate. Under the Black--Scholes--Merton benchmark,
\begin{equation}
 dS_t=(r-q)S_t\,dt+\sigma S_t\,dW_t^{Q}.
 \label{eq:q_gbm}
\end{equation}
The historical drift $\mu$ in Equation~\eqref{eq:p_gbm} does not enter Equation~\eqref{eq:q_gbm}. Posterior uncertainty about $\mu$ is therefore relevant for physical forecasting but is not propagated into the no-arbitrage price in the baseline model.

For a discretely monitored arithmetic-average option with fixing dates $t_1,\ldots,t_M$, conditional risk-neutral value is
\begin{equation}
 C^Q(\sigma)=e^{-r\tau}
 \mathbb E_{\sigma}^{Q}\!\left[\Pi(A_T,K)\mid\mathcal F_t\right],
 \label{eq:conditional_price}
\end{equation}
where $\Pi(A_T,K)$ is the call or put payoff and $\tau=T-t$. The posterior for $\sigma$ induces a posterior distribution over theoretical prices through the push-forward
\begin{equation}
 p(C\mid\mathcal D)=
 \int
 \delta_{C^Q(\sigma)}(C)
 p(\sigma\mid\mathcal D)\,d\sigma.
 \label{eq:pushforward}
\end{equation}
This distribution quantifies parameter uncertainty in a model price. It is not a predictive distribution of portfolio losses and is not interpreted as VaR or expected shortfall.

\subsection{WTI futures representation}

For the WTI APO application, the state variables are the first-nearby futures fixings in Equation~\eqref{eq:apo_average}. A parsimonious risk-neutral benchmark for a futures price is
\begin{equation}
 dF_u=\sigma F_u\,dW_u^Q.
 \label{eq:futures_q}
\end{equation}
The implementation does not replace the entire APO with one fixed CL contract. For each remaining fixing date it identifies the first-nearby contract applicable to that date and initializes the stochastic fixing from the corresponding point of the observed futures term structure.

For a valuation date inside the averaging month, simulation starts from Equation~\eqref{eq:partial_fixing}: realized fixings enter deterministically and only the remaining first-nearby settlements are stochastic. Two options with the same strike and final settlement date can therefore have different effective volatility exposure when their realized-average components differ.

The single-factor specification is deliberately interpretable. Under a constant-volatility GBM, the diffusion coefficient does not have to change mechanically when moving from $\mathbb P$ to $\mathbb Q$. Using a historical estimate of $\sigma$ in Equation~\eqref{eq:futures_q} is therefore coherent under the maintained model. It is not asserted to be a complete commodity model: mean reversion, stochastic convenience yield, stochastic rates, and jump or stochastic-volatility dynamics are well-established alternatives for commodity futures and Asian options \cite{schwartz1997,ewaldwu2023,shirayatakahashi2011}. A market-implied effective volatility can therefore differ because observed option prices reflect risk premia and because the one-factor constant-volatility benchmark may absorb omitted stochastic volatility, jumps, smile effects, or richer futures-curve dynamics. This distinction follows the broader option-pricing evidence that changing model specification can be a first-order empirical margin and that additional complexity must be judged by incremental performance rather than realism alone \cite{bakshicaochen1997,cumminsesposito2025}.

\subsection{Monte Carlo and conditioning-based valuation}

Arithmetic-average options generally lack the simple closed form available for geometric averages. The baseline empirical prices are therefore evaluated by Monte Carlo with common random numbers across nearby volatility values. In the standard Asian benchmark, antithetic sampling and a geometric-average control variate follow the logic of \cite{kemnavorst1990}. More generally, variance-reduction controls are used only when their expectation is known under the exact state representation being simulated.

For robustness, the same one-factor risk-neutral dynamics are also priced using the geometric-conditioning approximation of \cite{curran1994}. Curran pricing is not treated as the true exchange implementation; it is an independent numerical benchmark under the paper's maintained $Q$ dynamics. Difficult empirical contracts are further repriced using multi-million-path pseudo-random Monte Carlo and independently scrambled Sobol sequences. Agreement across these engines is used to determine whether the small posterior-integration effects are numerically resolved.

\subsection{Competing valuation rules}

Four baseline estimators are compared:
\begin{align}
 \widehat C_{PI}
 &=\mathbb E[C^Q(\sigma)\mid\mathcal D],\\
 \widehat C_{PM}
 &=C^Q(\mathbb E[\sigma\mid\mathcal D]),\\
 \widehat C_{Mode}
 &=C^Q(\widehat\sigma_{Mode}),\\
 \widehat C_{MLE}
 &=C^Q(\widehat\sigma_{MLE}).
 \label{eq:estimators}
\end{align}
The first is the posterior-integrated rule, sometimes called full-Bayes pricing. The second collapses the posterior to its mean before pricing. The third uses the marginal posterior mode of $\sigma$, and the fourth is a non-Bayesian likelihood benchmark. The marginal volatility mode is not described as a joint MAP estimator.

The main within-model diagnostic is
\begin{equation}
 \Delta_{PI,PM}
 =\widehat C_{PI}-\widehat C_{PM}.
 \label{eq:fb_pm_gap}
\end{equation}
A second-order expansion gives
\begin{equation}
 \Delta_{PI,PM}
 \approx
 \frac{1}{2}C^{Q\prime\prime}(\bar\sigma)
 \operatorname{Var}(\sigma\mid\mathcal D),
 \qquad
 \bar\sigma=\mathbb E[\sigma\mid\mathcal D].
 \label{eq:taylor_gap}
\end{equation}
Equation~\eqref{eq:taylor_gap} is the central mechanism: posterior integration matters only to the extent that posterior dispersion interacts with nonlinear exposure to volatility. The second derivative $C^{Q\prime\prime}(\sigma)$ is the option's volatility curvature, commonly termed \emph{volga} or \emph{vomma}; it need not have the same sign in every contract state.

This mean-price correction is distinct from the dispersion of the posterior distribution of model prices. A first-order delta-method approximation gives
\begin{equation}
 \operatorname{SD}\!\left(C^Q(\sigma)\mid\mathcal D\right)
 \approx
 \left|C^{Q\prime}(\bar\sigma)\right|
 \sqrt{\operatorname{Var}(\sigma\mid\mathcal D)}.
 \label{eq:price_uncertainty_delta}
\end{equation}
Thus the PI--PM correction is locally of order $\operatorname{Var}(\sigma\mid\mathcal D)$, whereas model-price uncertainty is locally of order $\sqrt{\operatorname{Var}(\sigma\mid\mathcal D)}$. Close PI and PM point prices therefore do not imply a narrow posterior distribution of theoretical prices.

As a plug-in sensitivity check, the paper also considers the posterior root-mean-square volatility
\begin{equation}
 \sigma_{RMS}
 =
 \sqrt{\mathbb E[\sigma^2\mid\mathcal D]}
 =
 \sqrt{\bar\sigma^2+\operatorname{Var}(\sigma\mid\mathcal D)},
 \label{eq:sigma_rms}
\end{equation}
with price $C^Q(\sigma_{RMS})$. This comparator is not a replacement for posterior integration; it tests how much the near-equivalence depends on using $\mathbb E[\sigma\mid\mathcal D]$ rather than a variance-matching scalar plug-in.

\subsection{APO-implied effective risk-neutral volatility}

To separate posterior integration from volatility specification, the paper also asks what volatility parameter the observed APO settlement requires under the same one-factor pricing map. For contract $i$ on date $t$, the contract-level implied value $\sigma^{APO}_{Q,i,t}$ solves
\begin{equation}
 C^Q_{i,t}(\sigma^{APO}_{Q,i,t})=P^{mkt}_{i,t},
 \label{eq:apo_iv}
\end{equation}
when the settlement lies within the monotone model price map. Cross-sectional common-volatility estimates minimize settlement squared error across contracts on a date, while smile specifications allow implied volatility to vary with moneyness and option type.

These quantities are best interpreted as \emph{Curran-equivalent APO-implied effective volatilities under the maintained one-factor model}. They are inferred from the APO cross-section itself and are therefore not an independent external $Q$-volatility source. Their role is diagnostic and predictive: they measure how far the volatility parameter needed by the option market lies from the historical constant-volatility estimate, and whether that information transfers to later valuation dates.

\subsection{Independent vanilla-WTI risk-neutral surface}\label{sec:external_lo_method}

The APO-implied exercise above is deliberately informative about volatility specification but remains within the target option family. To obtain an external risk-neutral state, the paper also uses standard monthly American WTI options on futures (CME product LO). For vanilla contract $j$ observed on trading reference date $s$, let $F_{j,s}$ be the matching CL futures settlement, $K_j$ the strike, $\tau_{j,s}$ time to expiration, and $r_s(\tau_{j,s})$ the date-specific rate proxy. The contract-level implied volatility $\sigma^{LO}_{Q,j,s}$ solves
\begin{equation}
 V^{A}_{LO}\!\left(F_{j,s},K_j,\tau_{j,s},r_s,\sigma^{LO}_{Q,j,s}\right)
 =P^{LO}_{j,s},
 \label{eq:lo_iv}
\end{equation}
where $V^{A}_{LO}$ is evaluated with an American Cox--Ross--Rubinstein futures-option tree and $P^{LO}_{j,s}$ is the official LO settlement. The production inversion uses 160 tree steps, as recorded in the inversion manifest. The futures process is a martingale under $\mathbb Q$ in the maintained one-factor representation, and early exercise is evaluated at every tree node.

The resulting contract-level IVs are fitted separately by underlying futures contract with the same quadratic log-moneyness design used for the APO forward smile, including call/put interactions. For an APO target on date $t$ and strike $K$, each required underlying $u$ therefore supplies a strictly prior implied volatility
\begin{equation}
 \widehat\sigma^{LO}_{Q,u,t-1}(K),
 \qquad s<t,
 \label{eq:lo_prior_guard}
\end{equation}
evaluated at moneyness computed from the target-date CL futures curve. The previous-day specification uses the latest completed LO date before $t$; the expanding specification uses all earlier dates with exponentially decaying recency weights.

The APO pricing engine retains one effective volatility state. If $w_u$ is the fraction of APO fixing dates mapped to underlying futures contract $u$, the strike-specific external state is
\begin{equation}
 \widehat\sigma^{LO,eff}_{Q,t}(K)
 =
 \left[
 \sum_u w_u
 \left\{\widehat\sigma^{LO}_{Q,u,t-1}(K)\right\}^2
 \right]^{1/2}.
 \label{eq:lo_effective_sigma}
\end{equation}
Equation~\eqref{eq:lo_effective_sigma} is an effective scalar transfer rule, not an exact variance identity for the arithmetic-average payoff. In general, if $U_t$ indexes unresolved fixings and $M$ is the total number of fixings,
\begin{equation}
 \operatorname{Var}_{Q}(A_T\mid\mathcal F_t)
 =
 \frac{1}{M^2}
 \sum_{j,k\in U_t}
 \operatorname{Cov}_{Q}\!\left(
 F_{t_j}^{u_j},F_{t_k}^{u_k}\mid\mathcal F_t
 \right),
 \label{eq:average_variance_covariance}
\end{equation}
so fixing times, futures levels, and cross-maturity dependence all matter. The weighted-RMS rule is used only to map maturity-specific vanilla IVs back into the paper's maintained one-factor scalar-volatility pricing engine. A weighted-mean aggregation is retained as a sensitivity specification.

The weights are derived from the target contract's actual fixing schedule rather than hard-coded. Predictions outside the prior observed LO moneyness support are clipped to the nearest support boundary and flagged; the main external comparison is repeated on the exact common-support subset with no clipping. No same-day APO implied volatility enters Equations~\eqref{eq:lo_iv}--\eqref{eq:lo_effective_sigma}.


\section{Research design}\label{sec:design}

\subsection{Synthetic repeated-sampling benchmark}

The simulation study first asks a question that cannot be answered cleanly with market prices alone: when the data-generating volatility is known, how much is lost by collapsing posterior uncertainty before pricing? Historical return samples are generated under $\mathbb P$ and option values are evaluated under $\mathbb Q$, preserving the measure separation in Section~\ref{sec:methodology}. The original production grid varies historical sample size, volatility, moneyness, and maturity and evaluates posterior-integrated, posterior-mean, posterior-mode, and MLE pricing against the external target
\begin{equation}
 C_0=C^Q(\sigma_0).
 \label{eq:synthetic_target}
\end{equation}
The target is not a posterior draw and is not constructed from any estimator being compared.

For method $M\in\{PI,PM,Mode,MLE\}$ and $R$ independent histories, the main statistics are
\begin{align}
 \operatorname{Bias}_M
 &=R^{-1}\sum_{j=1}^{R}(\widehat C_{M,j}-C_0),\\
 \operatorname{MAE}_M
 &=R^{-1}\sum_{j=1}^{R}|\widehat C_{M,j}-C_0|,\\
 \operatorname{RMSE}_M
 &=\left[R^{-1}\sum_{j=1}^{R}(\widehat C_{M,j}-C_0)^2\right]^{1/2}.
 \label{eq:metrics}
\end{align}

\subsection{Massive mechanism map}

The large mechanism experiment is designed around Equation~\eqref{eq:taylor_gap}. It uses
\begin{align}
 n&\in\{21,42,63,126,252,504,1260\},\\
 \sigma_0&\in\{0.10,0.20,0.35,0.50,0.80\},
\end{align}
with 5,000 independently generated histories in each $(n,\sigma_0)$ cell, for 175,000 synthetic datasets. The contract layer crosses four maturities (21, 63, 126, and 252 days), seven moneyness values from 0.60 to 1.50, and six partial-fixing states from zero to nearly the full averaging window.

Because the Gaussian prior on $\mu$ is conjugate to the likelihood conditional on $\sigma$, $\mu$ is integrated analytically. The remaining posterior in $\sigma$ is evaluated on a dense deterministic grid. This removes random-walk Metropolis error from the mechanism map while retaining exactly the baseline Gaussian statistical model. For each history and contract state the experiment computes $\Delta_{PI,PM}$, posterior variance of $\sigma$, local price curvature, and the Taylor approximation in Equation~\eqref{eq:taylor_gap}. The purpose is not merely to show that a gap can be made large, but to characterize the state variables that make posterior integration economically relevant.

\subsection{Empirical option panel}

The raw empirical sample contains every downloaded WTI APO history rather than a hand-selected set of strikes. The unit of observation is option $i$ on trading date $t$, identified by option type, strike, expiry month, and symbol. The raw panel contains 213 unique contracts and 11,179 option-date observations. The sample is intentionally unbalanced because strike availability and liquidity decline at longer maturities.

Three data layers remain separate. The \emph{raw panel} preserves all parsed observations. The \emph{main empirical panel} applies transparent data-quality and pricing-availability rules. The \emph{representative-contract set} is used only for figures and economic discussion. The October-2026 baseline contains 310 retained contract-date observations across 12 eligible valuation dates. Five dates contain at least one option with positive reported volume, and six individual contract-date observations report positive daily volume themselves.

\subsection{Reconstructing first-nearby fixings and moneyness}

Each option-date observation is linked to the information set required by the APO payoff. The pipeline reconstructs: (i) realized first-nearby settlements for fixing dates already elapsed, (ii) the relevant CL futures contracts for remaining fixing dates, and (iii) the contemporaneous futures term structure used to initialize those remaining fixings. This mapping produces $\widehat A^Q_{t,T}$ in Equation~\eqref{eq:expected_average} and the observation-specific moneyness measure in Equation~\eqref{eq:moneyness}.

The baseline physical-measure inference uses Yahoo \texttt{CL=F} as a labelled continuous/front-month return proxy; it is not used as a contractual substitute for the first-nearby sequence that determines the APO payoff. A dedicated source-robustness experiment reconstructs the historical first-nearby return series contract by contract from official final CL settlements. Each trading date is assigned to the earliest contract whose last-trade date has not passed. Returns spanning a contract switch are excluded so that contango or backwardation is not treated as a one-day diffusion shock. Four isolated mapped settlement gaps are retained in the audit trail but not imputed; any one-day return depending on a missing price is also excluded. The resulting 2024--2026 reconstruction contains 632 usable daily returns and 33 mapped roll switches, compared with 676 usable Yahoo returns over the same inference window.

Valuation uses committed Barchart individual CL histories with explicit last-trade dates. In the Barchart histories used here, \texttt{Latest} is the CME settlement field; the source-field name is retained for provenance and is not interpreted as an intraday last trade.

\subsection{Rolling physical-measure uncertainty}

Let $\mathcal D_t$ denote the historical information set available at date $t$. For an expanding sample beginning in January 2024, the posterior
\begin{equation}
 p(\sigma\mid\mathcal D_t)
\end{equation}
is estimated without using future returns. Each option-date observation is then priced using posterior integration, posterior-mean volatility, the marginal posterior mode, and the historical MLE. This rolling construction prevents later returns from contaminating earlier valuation dates.

The principal pricing errors are
\begin{equation}
 e^M_{i,t}=P^{M}_{i,t}-P^{mkt}_{i,t},
\end{equation}
where $P^{mkt}_{i,t}$ is the Barchart end-of-day settlement field and $M$ indexes the pricing rule. Mean error, MAE, and RMSE are reported overall and within liquidity-restricted samples.

To separate information source from information recency, a pre-specified recent-history benchmark reuses the exact strict-forward holdouts and target-date information cutoffs. It evaluates annualized realized volatility from the most recent 63, 126, and 252 usable daily returns, plus an exponentially weighted estimate with a 63-business-day half-life. These choices are fixed before inspecting their pricing errors. Each scalar estimate is mapped through the same Curran APO pricing engine used for the prior-date smile predictions. The purpose is not to optimize a historical window ex post, but to ask whether a reasonable recent historical benchmark materially narrows the option-informed advantage.

A second scalar sensitivity uses Equation~\eqref{eq:sigma_rms}. Because every empirical run already stores the posterior mean and standard deviation of $\sigma$, the posterior-RMS comparator can be evaluated on the same holdouts without rerunning MCMC.

\subsection{APO-implied volatility and cross-sectional validation}

The initial September and October panels provide 402 contract-date observations for implied-volatility analysis. Contract-level implied volatilities are obtained by inverting Equation~\eqref{eq:apo_iv}. A date-level common $\sigma_Q$ is then calibrated by minimizing cross-sectional settlement squared error. Two deliberately harder checks avoid evaluating a contract on information extracted from that same contract: leave-one-contract-out calibration excludes the target option, and cross-option-type transfer estimates from calls and evaluates puts or vice versa.

These exercises remain cross-sectional. The primary predictive robustness design is therefore strictly forward in time. For each target date and expiry, the volatility smile is estimated using only APO implied-volatility observations from dates strictly earlier than the target. Two schemes are considered: a \emph{previous-day smile} using the latest earlier completed date and an \emph{expanding smile} using all earlier dates with exponential recency weights. Same-day implied volatility is retained only as an ex-post diagnostic and never enters the target-date fit.

A persistence benchmark asks whether the fitted previous-day smile adds information beyond simply carrying forward the same contract's latest strictly prior APO implied volatility. The carried volatility is repriced using the target-date futures curve, realized-fixing state, discount factor, and remaining fixing schedule. Thus the benchmark preserves a transparent current-state adjustment while excluding same-day APO information. A committed-data coverage audit requires at least six main-sample contracts on a date before an expiry is admitted to the smile experiment. Seven expiries pass that screen---September, October, and November 2026; March and September 2027; and March and September 2028---while June 2029 does not. The extended evaluation contains 2,381 option-date holdouts over 15 distinct target dates, with results reported both pooled and separately by expiry.

\subsection{Independent vanilla-WTI forward validation}\label{sec:external_lo_design}

The APO prior-date smile establishes that option-market volatility information is forward-relevant, but it does not by itself rule out a within-family explanation. The independent validation therefore constructs a risk-neutral state from standard monthly WTI futures options (LO) and uses it to predict October-2026 APO settlements. Raw vendor records are kept local; only reproducible query metadata, diagnostics, and derived results are versioned.

The acquisition window is August 24 through September 10, 2026, with CME trading reference date used for temporal ordering. The final strike range is 70--120 on the two CL futures contracts required by the October APO fixing schedule, CLX6 and CLZ6. Final official LO settlements are paired with official CL futures settlements and inverted using Equation~\eqref{eq:lo_iv}. The resulting panel contains 5,220 successful contract-date implied-volatility inversions over 13 reference dates. Of these settlements, 5,214 are flagged actual and six theoretical.

For each target APO date $t$, every LO observation used to fit the external surface satisfies $s<t$. The \emph{previous-day} surface uses only the latest completed LO reference date; the \emph{expanding} surface uses all earlier reference dates with a five-calendar-day exponential half-life. The primary specification evaluates the fitted surface separately for CLX6 and CLZ6 at each APO strike and combines the component volatilities with the fixing-count RMS rule in Equation~\eqref{eq:lo_effective_sigma}. A pre-specified sensitivity instead chooses the scalar volatility that matches the variance of the unresolved arithmetic-average component under the maintained common-factor representation, using target-date futures levels and fixing times. Same-day APO prices are never used to estimate the external state.

The final October comparison is matched at the contract-date level. Across the 280 observations shared by the historical baseline, the two external LO specifications, and the two prior-date APO smiles, 253 contract-dates lie inside the observed prior LO moneyness support for both external surface specifications. This 90.4\% no-clipping subset spans all 10 matched valuation dates and is the primary external-validation sample. Error differences are resampled by valuation-date cluster with 100,000 bootstrap replications so that contracts sharing one market date are not treated as independent observations.

\subsection{Heavy-tail robustness under the physical measure}

To test whether the Gaussian historical likelihood drives the pricing hierarchy, the physical-measure innovation is replaced by a standardized Student-$t$ variable,
\begin{equation}
 R_i=
 \left(\mu-\frac12\sigma^2\right)\Delta t
 +\sigma\sqrt{\Delta t}\,\varepsilon_i,
 \qquad
 \operatorname{Var}(\varepsilon_i)=1,
\end{equation}
with degrees of freedom inferred jointly with $\mu$ and $\sigma$. The $Q$ pricing dynamics, realized fixings, futures curve, and pricing map are otherwise unchanged. The production robustness run uses eight chains of 100,000 iterations per unique valuation date and covers 402 option-date observations across the September and October expiries.

\subsection{Numerical accuracy and posterior calibration}

Small values of $\Delta_{PI,PM}$ are scientifically useful only if they are larger than pricing noise. Eighteen empirically difficult contracts are therefore repriced over a dense volatility grid with independent high-precision pseudo-random Monte Carlo runs. Eight common difficult targets are also evaluated with independently scrambled Sobol sequences, and Curran conditioning is evaluated on the same volatility grids. The budgets are chosen so that direct Monte Carlo uncertainty in the PI--PM gap is typically of order $10^{-6}$, well below the observed empirical gaps near $10^{-3}$.

Posterior implementation is checked separately by simulation-based calibration. Two hundred thousand datasets are drawn from the prior predictive distribution, with 50,000 replications at each of $n=21,63,252,1260$. Calibration is assessed through posterior-CDF ranks and 50\%, 80\%, and 95\% credible-interval coverage.

\subsection{Liquidity dependence and clustered uncertainty}

Contracts observed on the same valuation date share the futures curve, historical posterior, and market shock, so treating option rows as independent would overstate precision. Error comparisons are therefore bootstrapped by valuation-date cluster rather than by individual contract. The production robustness analysis uses 100,000 cluster resamples for the full sample, the positive-volume subset, and open-interest thresholds of 1, 10, 100, and 500 contracts. The bootstrap reports percentile intervals and the probability that a candidate specification improves MAE or RMSE relative to the historical-volatility posterior-integrated baseline.

\section{Results}\label{sec:results}

The completed evidence separates four questions. First, when does posterior integration differ materially from posterior-mean plug-in pricing? Second, can a small PI--PM mean correction coexist with substantial posterior model-price dispersion? Third, how large is the integration effect relative to changing the volatility input, including stronger recent-history and persistence benchmarks? Fourth, does the option-informed result survive when the volatility state comes from an independent vanilla-WTI market?

\subsection{Repeated-sampling benchmark}

Table~\ref{tab:synthetic_pilot} reports the original repeated-sampling benchmark. Even with only 63 daily observations, posterior-integrated and posterior-mean plug-in pricing are nearly indistinguishable. As the historical sample increases, all four estimators converge and the PI--PM difference becomes numerically negligible.

\begin{table*}[t]
\centering
\caption{Repeated-sampling point-pricing performance. The target is the risk-neutral price evaluated at the true volatility, not a posterior-generated target. The sigma-mode comparator is the marginal posterior mode of volatility rather than a joint MAP estimate. Posterior interval coverage is reported separately in the text because the interval procedure is common across point estimators.}\label{tab:synthetic_pilot}
\begin{tabular}{llrrr}
\toprule
$N$ & Method & Bias & MAE & RMSE\\
\midrule
63 & Posterior-integrated & -0.0011 & 0.3543 & 0.4686\\
63 & Posterior-mean plug-in & -0.0012 & 0.3543 & 0.4687\\
63 & Sigma posterior-mode plug-in & -0.0919 & 0.3783 & 0.4867\\
63 & MLE plug-in & -0.0407 & 0.3576 & 0.4701\\
\addlinespace
252 & Posterior-integrated & 0.0139 & 0.1988 & 0.2454\\
252 & Posterior-mean plug-in & 0.0139 & 0.1988 & 0.2454\\
252 & Sigma posterior-mode plug-in & -0.0212 & 0.1974 & 0.2446\\
252 & MLE plug-in & 0.0035 & 0.1992 & 0.2444\\
\addlinespace
1260 & Posterior-integrated & -0.0048 & 0.1010 & 0.1243\\
1260 & Posterior-mean plug-in & -0.0048 & 0.1010 & 0.1243\\
1260 & Sigma posterior-mode plug-in & -0.0128 & 0.1060 & 0.1302\\
1260 & MLE plug-in & -0.0067 & 0.1013 & 0.1243\\
\bottomrule
\end{tabular}
\end{table*}

The central 95\% posterior model-price interval has empirical coverage 0.9500, 0.9625, and 0.9250 for $N=63,252,$ and $1260$, respectively. These are properties of the posterior interval procedure and are not separate coverage statistics for each point estimator.

The near equality of $\widehat C_{PI}$ and $\widehat C_{PM}$ is not a generic Bayesian result. It is a property of the posterior dispersion and pricing curvature in these particular states. The massive mechanism experiment makes this distinction explicit.

\subsection{Where posterior integration matters}

The mechanism map covers 175,000 independently simulated historical datasets and 168 WTI-like contract states. Across the complete experiment, the largest observed absolute PI--PM gap is 0.1384. The largest cell-level 99th percentile is 0.1271. These values are orders of magnitude larger than the roughly $10^{-3}$ gaps observed in the WTI baseline, so the empirical near-equivalence of PI and PM is not an artifact of a pricing rule that can never separate.

The mechanism map normalizes the WTI-like forward level to \$100 per barrel, so the maximum gap of 0.1384 is \$0.1384 per barrel. Under the WTI APO contract multiplier of 1,000 barrels and ordinary minimum price fluctuation of \$0.01 per barrel, that scale corresponds to approximately \$138.40 per contract, or 13.84 ordinary ticks \cite{cme_wti_apo_rule}. This conversion gives the synthetic separation an explicit economic scale; it does not imply that every state attaining the statistical maximum would be equally tradable.

The adverse states have a clear economic structure. The largest gaps occur with short historical samples, high true volatility, long remaining horizons, little realized fixing, and strikes far from the expected average. For example, with $n=21$, $\sigma_0=0.80$, 126 days to maturity, no realized fixing, and moneyness 1.50, the mean PI--PM gap is 0.1176 and the 99th percentile is 0.1271. With $n=21$, $\sigma_0=0.50$, 252 days, no realized fixing, and moneyness 1.50, the mean gap is 0.1103 and the maximum observed gap is 0.1384.

The mechanism in Equation~\eqref{eq:taylor_gap} explains these states remarkably well. In the most adverse cells, the correlation between the realized PI--PM gap and the second-order approximation is approximately 0.999, with through-origin slopes close to one. Posterior dispersion also behaves as expected: for $\sigma_0=0.80$, mean posterior standard deviation falls from approximately 0.122 at $n=21$ to 0.0159 at $n=1260$. The experiment therefore turns the qualitative Jensen argument into a quantitative state map: integration matters when \emph{both} posterior dispersion and price curvature are large.

Partial fixing works in the opposite direction. As a larger fraction of the final arithmetic average becomes deterministic, the remaining volatility exposure shrinks and so does the opportunity for posterior integration to change the point price. This effect is visible even in high-volatility, extreme-moneyness cells and provides a contract-specific reason why parameter uncertainty can matter early in an averaging problem but fade rapidly as fixings accumulate.

Figure~\ref{fig:mechanism_map} condenses these mechanism results. Panel A shows posterior dispersion shrinking with historical information; Panels B and C show how the integration gap grows in high-curvature, weakly fixed states and then contracts as fixings accumulate; Panel D compares the realized PI--PM gap with the second-order approximation across the complete synthetic state map.

\begin{figure*}[t]
\centering
\includegraphics[width=\textwidth]{figures/publication/fig01_mechanism_map.pdf}
\caption{Mechanism map for posterior integration. Panel A reports the mean posterior standard deviation of $\sigma$ against historical sample size for the simulated volatility levels. Panel B shows mean $|PI-PM|$ across moneyness and partial-fixing states for $n=21$, $\sigma_0=0.80$, and 126 days to maturity. Panel C traces the same integration gap against the fraction already fixed at $K/F=1.50$ for the four maturity horizons. Panel D compares the actual cell-level mean PI--PM gap with the second-order Taylor term in Equation~\eqref{eq:taylor_gap}; the correlation across synthetic cells is 0.9995. The figure uses only independently simulated histories and prices under the maintained one-factor risk-neutral map.}
\label{fig:mechanism_map}
\end{figure*}

\subsection{Contract-consistent October-2026 baseline}

The first completed WTI baseline uses the October-2026 APO expiry. Across 12 eligible valuation dates, the main sample contains 310 option-date observations. Five dates contain at least one positive-volume contract, producing a 136-observation positive-volume-date sample, while only six individual contract-date observations themselves report positive daily volume.

Table~\ref{tab:oct26_empirical} reports pooled model-to-market errors. All rules use the same futures curves, realized fixings, discount factors, and market settlements; the comparison therefore isolates the incremental effect of the volatility decision rule within the maintained model.

\begin{table*}[t]
\centering
\caption{October-2026 WTI APO point-pricing errors and posterior price uncertainty. Mean error is model price minus the Barchart end-of-day settlement field. Panel B compares the nonlinear PI--PM mean-price correction with the width of the central 95\% posterior distribution of theoretical prices.}\label{tab:oct26_empirical}

\textit{Panel A: point-pricing errors}\\[2pt]
\begin{tabular}{llrrrr}
\toprule
Sample & Method & $N$ & Mean error & MAE & RMSE\\
\midrule
All eligible dates
 & Posterior-integrated & 310 & -0.4065 & 0.4079 & 0.5177\\
 & Posterior mean & 310 & -0.4072 & 0.4086 & 0.5183\\
 & Sigma posterior mode & 310 & -0.4122 & 0.4132 & 0.5200\\
 & MLE & 310 & -0.4088 & 0.4101 & 0.5197\\
\addlinespace
Positive-volume dates
 & Posterior-integrated & 136 & -0.3601 & 0.3633 & 0.4303\\
 & Posterior mean & 136 & -0.3609 & 0.3640 & 0.4310\\
 & Sigma posterior mode & 136 & -0.3697 & 0.3720 & 0.4413\\
 & MLE & 136 & -0.3625 & 0.3654 & 0.4324\\
\addlinespace
Positive-volume contracts
 & Posterior-integrated & 6 & -0.1388 & 0.1543 & 0.1864\\
 & Posterior mean & 6 & -0.1394 & 0.1547 & 0.1869\\
 & Sigma posterior mode & 6 & -0.1530 & 0.1651 & 0.2026\\
 & MLE & 6 & -0.1414 & 0.1560 & 0.1882\\
\bottomrule
\end{tabular}

\vspace{5pt}
\textit{Panel B: mean-price correction versus posterior price dispersion}\\[2pt]
\begin{tabular}{lrrrr}
\toprule
Sample & $N$ & Mean $|PI-PM|$ & Mean 95\% width & Width / gap\\
\midrule
All eligible dates & 310 & 0.000745 & 0.2733 & 366.7\\
Positive-volume dates & 136 & 0.000802 & 0.2695 & 336.1\\
Positive-volume contracts & 6 & 0.000591 & 0.3660 & 619.7\\
\bottomrule
\end{tabular}
\end{table*}

Posterior integration and the posterior-mean plug-in are extremely close in every sample, but Panel B shows why that fact must not be read as absence of parameter uncertainty. Across all 310 observations, the mean central 95\% posterior price-interval width is 0.2733, roughly 367 times the mean absolute PI--PM correction. Among the six positive-volume contracts, the corresponding ratio is about 620. The empirical conclusion is therefore specific: the nonlinear shift in the posterior \emph{mean} price is small under the observed posterior and contract states, while the posterior distribution of theoretical model prices remains visibly dispersed.

The sign of the model-to-market errors is equally important. All four historical-volatility rules underprice the option settlements on average. This common bias cannot be explained by choosing PI rather than PM because the two rules produce almost the same prices. It points instead toward a larger specification discrepancy.

Because the empirical PI--PM differences are sub-cent, Section~\ref{sec:numerical_robustness} subjects them to high-precision pseudo-random Monte Carlo, randomized Sobol, and Curran benchmarks. The three engines recover the same integration gap to a scale far smaller than the gap itself, so the near-equivalence is not a simulation-noise artifact.

\subsection{Historical versus APO-implied effective volatility}

The September--October APO cross-section provides 402 contract-date observations for implied-volatility inversion. All 402 settlements can be mapped to a finite implied volatility under the maintained Curran-equivalent one-factor model. The median historical posterior mean is 0.4155. By contrast, the median date-level common volatility calibrated to the complete APO cross-section is 0.4577, a median difference of 0.0425. Contract-level implied volatilities have a still higher median of 0.4835 because the cross section contains pronounced strike variation.

These values should not be interpreted as proof that the diffusion coefficient must differ between $\mathbb P$ and $\mathbb Q$. They are effective parameters required by the maintained pricing map. Their economic importance is that they expose a specification margin much larger than the PI--PM correction.

Figure~\ref{fig:historical_vs_apo_iv} shows this distinction at both the date and strike levels. The historical posterior mean remains comparatively stable, whereas the model-equivalent APO-implied state moves with valuation date and exhibits clear strike variation.

\begin{figure*}[t]
\centering
\includegraphics[width=\textwidth]{figures/publication/fig02_historical_vs_apo_implied_volatility.pdf}
\caption{Historical and APO-implied effective volatility in the September--October 2026 cross-section. Panels A and B compare the historical posterior mean $\sigma_P$ with the date-level common APO-implied $\sigma_Q$ for the September and October expiries. Panels C and D show contract-level APO-implied volatilities against log moneyness on August 26 and September 10, with horizontal references for the historical posterior mean and the date-level common APO-implied value. These implied volatilities are effective parameters under the maintained Curran-equivalent pricing map; they are not interpreted as unique structural diffusion coefficients.}
\label{fig:historical_vs_apo_iv}
\end{figure*}

Leave-one-contract-out calibration provides a first cross-sectional check. Across 399 targets, a common APO-calibrated volatility reduces MAE from 0.3450 for the historical-volatility PI baseline to 0.1557 and RMSE from 0.4625 to 0.1900. The improvement is not uniform, however. The nine positive-volume observations move in the opposite direction in aggregate, and cross-option-type transfer is weaker. These failures are consistent with the fact that one scalar volatility cannot reproduce the observed strike structure and that call/put samples occupy different moneyness regions.

\subsection{Strict forward-in-time volatility validation}

The stronger test predicts each target date using only implied-volatility observations from earlier dates. The extended experiment spans seven APO expiries from September 2026 through September 2028 and contains 2,381 holdout option-date observations. No same-day option settlement enters the fitted smile.

The extension reveals a pronounced term structure in the effective volatility required by the maintained pricing map. Median date-level common $\sigma_Q$ is 0.4478, 0.4572, and 0.4444 for the September, October, and November 2026 expiries, respectively; it then declines to 0.3609 for March 2027, 0.3001 for September 2027, 0.2647 for March 2028, and 0.2416 for September 2028. Over the same valuation dates, the historical posterior mean remains near 0.416. These $\sigma_Q$ values are model-equivalent parameters, not structural volatility estimates; the maturity pattern can absorb risk premia, smile effects, omitted futures-curve dynamics, and liquidity or marking effects.

Table~\ref{tab:forward_q} reports the corresponding strict-forward pricing errors by expiry.

\begin{table*}[t]
\centering
\caption{Strict forward-in-time APO volatility validation by expiry. The historical-volatility baseline is the posterior-integrated price on the same holdout observations.}\label{tab:forward_q}
\begin{tabular}{lrrrrrrr}
\toprule
Expiry & $N$ & \multicolumn{3}{c}{MAE} & \multicolumn{3}{c}{RMSE}\\
\cmidrule(lr){3-5}\cmidrule(lr){6-8}
 & & Historical PI & Expanding & Previous day & Historical PI & Expanding & Previous day\\
\midrule
2026-09 & 198 & 0.1090 & 0.0433 & 0.0443 & 0.1378 & 0.0754 & 0.0695\\
2026-10 & 280 & 0.4159 & 0.1486 & 0.1543 & 0.5308 & 0.2772 & 0.2557\\
2026-11 & 285 & 0.4276 & 0.1653 & 0.1729 & 0.5267 & 0.2696 & 0.2548\\
2027-03 & 339 & 1.0599 & 0.1187 & 0.1070 & 1.0744 & 0.1560 & 0.1333\\
2027-09 & 429 & 2.8963 & 0.0789 & 0.0867 & 2.9324 & 0.0979 & 0.1012\\
2028-03 & 445 & 4.5398 & 0.1096 & 0.1022 & 4.5813 & 0.1392 & 0.1248\\
2028-09 & 405 & 5.8103 & 0.0962 & 0.0885 & 5.8979 & 0.1268 & 0.1152\\
\addlinespace
Pooled & 2,381 & 2.6187 & 0.1088 & 0.1075 & 3.4090 & 0.1725 & 0.1594\\
\bottomrule
\end{tabular}
\end{table*}

Figure~\ref{fig:forward_q_bootstrap} complements the point estimates with valuation-date cluster uncertainty. Negative differences indicate lower error for the prior-date option-implied specification than for historical-volatility posterior integration.

\begin{figure*}[t]
\centering
\includegraphics[width=\textwidth]{figures/publication/fig03_forward_q_cluster_bootstrap.pdf}
\caption{Strict forward-in-time pricing improvement across the seven admitted APO expiries. Points report the difference between each prior-date APO-smile error and the historical posterior-integrated baseline on the same holdouts; Panel A uses MAE and Panel B RMSE. Horizontal bars are 95\% percentile intervals from 100,000 valuation-date cluster bootstrap resamples, preserving dependence among contracts observed on the same market date. The dotted vertical line at zero is the no-difference benchmark: values to its left indicate lower error for the option-informed specification, whereas values to its right indicate worse performance relative to the historical-volatility baseline. Both the expanding and previous-day specifications use only implied-volatility observations dated strictly before each target date.}
\label{fig:forward_q_bootstrap}
\end{figure*}

The result is not driven by one expiry. Both prior-date specifications improve MAE and RMSE relative to the historical-volatility baseline in every admitted expiry. In the pooled panel, the expanding smile lowers MAE by 95.8\% and RMSE by 94.9\%, while the previous-day smile lowers MAE by 95.9\% and RMSE by 95.3\%. Mean pricing error moves from +2.4030 under the historical-volatility baseline to $-0.0697$ and $-0.0352$, respectively. The sign change relative to the short-dated October baseline reflects the maturity structure above: a long-window historical volatility near 0.416 is below the effective APO level for some 2026 expiries but substantially above the effective level embedded in the 2027--2028 marks.

A stronger benchmark asks whether this difference is merely a recency effect. Table~\ref{tab:recent_history_persistence} Panel A evaluates four pre-specified historical estimators on the exact same 2,381 holdouts: realized volatility from the latest 63, 126, and 252 usable returns and an exponentially weighted estimate with a 63-business-day half-life. None closes the gap. The best pooled recent-history rule is the 63-return estimator, with MAE 4.5201 and RMSE 5.6465, both worse than the long-window historical PI benchmark. The recent estimates are themselves high---their pooled mean volatilities range from 0.5141 to 0.6962---so shortening the history tends to aggravate the overpricing of the 2027--2028 expiries rather than reproduce the lower effective volatility embedded in those marks. Thus the forward advantage of the option-informed inputs cannot be attributed simply to using more recent data.

\begin{table*}[t]
\centering
\caption{Recent-history and persistence benchmarks for the strict-forward APO experiment. Panel A uses the same 2,381 contract-date holdouts and 15 target dates as the main forward comparison. Panel B restricts all methods to the 2,366 contract-dates for which the same contract has a strictly prior implied volatility. All option-informed inputs are dated strictly before the target date.}\label{tab:recent_history_persistence}

\textit{Panel A: recent historical volatility on the exact forward holdouts}\\[2pt]
\begin{tabular}{lrrrr}
\toprule
Volatility input & $N$ & Dates & MAE & RMSE\\
\midrule
Long-window historical PI & 2,381 & 15 & 2.6187 & 3.4090\\
Rolling 252 returns & 2,381 & 15 & 5.2183 & 6.3933\\
Rolling 126 returns & 2,381 & 15 & 8.3419 & 9.7846\\
Rolling 63 returns & 2,381 & 15 & 4.5201 & 5.6465\\
EWMA, half-life 63 days & 2,381 & 15 & 5.6616 & 6.8824\\
Prior-date APO expanding smile & 2,381 & 15 & 0.1088 & 0.1725\\
Prior-date APO previous-day smile & 2,381 & 15 & 0.1075 & 0.1594\\
\bottomrule
\end{tabular}

\vspace{5pt}
\textit{Panel B: contract-level persistence on exact common rows}\\[2pt]
\begin{tabular}{lrrrr}
\toprule
Volatility input & $N$ & Dates & MAE & RMSE\\
\midrule
Historical PI & 2,366 & 15 & 2.6279 & 3.4173\\
Prior-contract IV carry & 2,366 & 15 & 0.0992 & 0.1403\\
Previous-day fitted APO smile & 2,366 & 15 & 0.1074 & 0.1594\\
\bottomrule
\end{tabular}
\end{table*}

Panel B separates option information from cross-sectional fitting flexibility. A same-contract persistence rule carries the latest strictly prior APO implied volatility forward and reprices the contract using the target-date futures curve, discount factor, realized fixings, and remaining fixing schedule. It is available for 2,366 of the 2,381 holdouts. On those exact common rows, prior-contract IV carry has MAE 0.0992 and RMSE 0.1403, compared with 0.1074 and 0.1594 for the fitted previous-day smile. The fitted smile therefore does not drive the option-informed advantage; if anything, the simpler contract-level persistence rule is slightly more accurate in the pooled sample. The empirical content is consequently narrower and stronger: recent option-implied states are highly persistent and informative for subsequent APO marks, while flexible cross-sectional smoothing is not required to obtain the large improvement over historical-volatility inputs.

Section~\ref{sec:robustness} reconstructs the physical return history contract by contract from official first-nearby CL settlements. On these exact 2,381 holdouts, the reconstructed historical baseline has MAE 2.6468 and RMSE 3.4451 rather than 2.6187 and 3.4090. The option-informed errors are unchanged, so the central forward comparison is insensitive to the continuous-contract return proxy.

Liquidity restrictions preserve the central result. The 62 positive-volume holdouts span 13 valuation-date clusters: historical PI has MAE 1.6340 and RMSE 2.3391, compared with 0.1003 and 0.1358 for the expanding smile and 0.1084 and 0.1333 for the previous-day smile. At open interest of at least 100, the corresponding forward MAEs are approximately 0.106--0.107 versus 1.990 for the historical baseline. The long-dated positive-volume counts remain sparse within individual expiries, however, so the strongest statement concerns the pooled transaction-active sample and the much denser open-interest-filtered cross sections.

The forward result changes the interpretation of the empirical exercise. The observed pricing errors are not primarily a consequence of failing to integrate the historical volatility posterior, and the recent-history experiment shows that recency alone does not account for the difference. Prior-date option-implied states explain a much larger share of the subsequent pricing discrepancy across the observed maturity range. The persistence benchmark further shows that this result is not created by the flexibility of the fitted smile: carrying a contract's own prior implied volatility performs at least as well in the pooled common sample.

This comparison does not eliminate a role for Bayesian parameter uncertainty. The mechanism map shows where the posterior-integration correction can become appreciable. In the observed holdouts, however, changing from historical volatility inputs to recent option-implied states changes pricing errors far more than integrating the relatively concentrated historical posterior through the local pricing map. This remains an empirical benchmark comparison rather than a pure structural decomposition, because option-implied and historical inputs differ in information source and may embed risk premia, omitted dynamics, and market-marking effects.

\subsection{Independent vanilla-WTI validation}\label{sec:external_lo_results}

The strict APO-forward experiment leaves one important identification concern: its risk-neutral state is inferred from the same option family later used as the pricing target. The vanilla-WTI experiment addresses that concern by constructing the volatility surface only from prior-date LO and CL settlements. The final strike expansion to 70--120 places 253 of the 280 matched October contract-dates inside the observed prior LO moneyness support, compared with only 60 observations under the initial narrow strike pilot. Only 50 of 620 external-surface prediction rows require any support clipping.

Table~\ref{tab:external_lo} reports the exact no-clipping comparison. Historical-volatility posterior integration has MAE/RMSE 0.4464/0.5562 with Yahoo \texttt{CL=F} and 0.4362/0.5466 with the reconstructed first-nearby history. The external previous-day LO surface lowers these errors to 0.1559/0.2599, while the expanding LO surface gives 0.1604/0.2858. The corresponding prior-date APO smiles are 0.1676/0.2686 and 0.1612/0.2912. Thus an independently observed vanilla-WTI risk-neutral surface reproduces the large pricing improvement previously obtained from the APO market itself, and that comparison is insensitive to the historical return-source construction.

\begin{table*}[t]
\centering
\caption{Independent vanilla-WTI validation on the exact common-support sample. All methods are evaluated on the same 253 October-2026 APO contract-dates over 10 valuation dates. The two historical rows differ only in the physical-return source: Yahoo \texttt{CL=F} versus the contract-reconstructed first-nearby CL settlement series. External LO surfaces use only vanilla observations dated strictly before the APO target date.}\label{tab:external_lo}
\begin{tabular}{lrrr}
\toprule
Method & $N$ & MAE & RMSE\\
\midrule
Historical PI, Yahoo \texttt{CL=F} & 253 & 0.4464 & 0.5562\\
Historical PI, reconstructed first-nearby & 253 & 0.4362 & 0.5466\\
External LO surface, previous day & 253 & 0.1559 & 0.2599\\
External LO surface, expanding & 253 & 0.1604 & 0.2858\\
Prior-date APO smile, previous day & 253 & 0.1676 & 0.2686\\
Prior-date APO smile, expanding & 253 & 0.1612 & 0.2912\\
\bottomrule
\end{tabular}
\end{table*}

The valuation-date cluster bootstrap confirms that the external improvement over historical volatility is not driven by treating same-day contracts as independent. For the previous-day LO surface, the MAE difference relative to historical PI is $-0.2905$, with a 95\% interval of $[-0.3439,-0.2333]$, and the RMSE difference is $-0.2963$, with interval $[-0.3503,-0.2580]$. The expanding LO surface produces corresponding differences of $-0.2860$ and $-0.2704$, again with intervals entirely below zero.

The external and APO-implied surfaces are much closer to one another. Previous-day LO has an MAE difference of $-0.0117$ relative to the previous-day APO smile, with a 95\% interval of $[-0.0251,-0.0001]$; its RMSE difference is $-0.0087$, with interval $[-0.0149,0.0002]$. For the expanding specifications, both the MAE and RMSE intervals include zero. These results do not establish global superiority of vanilla LO over APO-implied volatility. Their more important implication is that the pricing gain associated with an option-informed risk-neutral state survives when the volatility source is moved outside the APO family.

The independent LO evidence leads to the same narrower comparison: the point-price effect of integrating the maintained historical-volatility posterior is much smaller than the pricing difference obtained from recent option-informed volatility inputs. This statement is supported by both within-family forward APO evidence and cross-instrument LO evidence. Bayesian integration propagates uncertainty conditional on the maintained model; it does not by itself resolve a volatility input that is empirically misaligned with option settlements.


\section{Robustness and extensions}\label{sec:robustness}

The completed robustness analysis asks whether the paper's main hierarchy survives changes in statistical assumptions, numerical pricing engines, posterior calibration checks, and liquidity weighting. The purpose is not to search for a specification that mechanically minimizes in-sample error, but to determine which conclusions are stable and which remain sample-dependent.

\subsection{Prior and estimation-window sensitivity}

The inverse-gamma prior in Equation~\eqref{eq:priors} is not treated as uniquely correct. A dedicated sensitivity experiment compares the baseline prior with two alternatives centered near 0.40 and estimation windows of 63, 126, 252, 504, and 673 daily returns. Four chains of 20,000 iterations are used for each specification, with maximum observed $\hat R_\sigma$ of 1.0007.

The main distinction is between prior sensitivity and window sensitivity. With 673 returns, posterior mean volatility is 0.4151 under the baseline prior, 0.4153 under a broad WTI-centered prior, and 0.4153 under a more concentrated WTI-centered prior. The corresponding Curran posterior-integrated prices are 0.6419, 0.6430, and 0.6428. Thus reasonable prior changes are negligible once the likelihood is informative. With only 63 observations, prior effects become visible at the cent level, while the PI--PM gap itself reaches approximately 0.0146. Changes in the historical window are materially larger because the windows span different realized volatility regimes.

This experiment prices one representative difficult contract and therefore diagnoses sensitivity of the posterior-integration mechanism; it is not the same as a comprehensive forward comparison on all 2,381 holdouts. The exact-holdout recent-history benchmark in Section~\ref{sec:design} is designed to address that distinct question without selecting the historical window after seeing pricing errors.

\subsection{Heavy-tailed physical-measure returns}

WTI returns are not especially well described by Gaussian innovations. Replacing the Gaussian likelihood under $\mathbb P$ with standardized Student-$t$ innovations produces posterior mean degrees of freedom around 3.8 across valuation dates, with posterior medians near 3.74--3.77. The data therefore support substantially heavier tails than the Gaussian benchmark.

The change does not overturn the main pricing hierarchy. Across 402 September--October contract-date observations, Student-$t$ posterior-integrated pricing has MAE 0.3263 and RMSE 0.4418, compared with 0.3444 and 0.4616 for the Gaussian posterior-integrated baseline. The improvement is approximately 5\% in MAE and 4\% in RMSE. Among nine positive-volume observations, MAE falls from 0.1249 to 0.1125 and RMSE from 0.1572 to 0.1413. Heavy tails therefore improve the physical-measure volatility estimate modestly, but the remaining error is still far larger than the within-model PI--PM correction.

The Student-$t$ volatility posterior is also wider: posterior standard deviations are roughly 0.027 rather than approximately 0.011 in the Gaussian long-window baseline. Even with that additional uncertainty, posterior-integrated and posterior-mean Student-$t$ prices remain close. This provides a direct robustness check on the claim that the observed pricing map is locally smooth relative to posterior dispersion.

\subsection{Contract-reconstructed physical volatility}

The baseline physical-measure analysis uses Yahoo \texttt{CL=F} as a reproducible continuous/front-month proxy, so a natural concern is that its undocumented roll construction may distort the historical volatility benchmark. A dedicated robustness experiment reconstructs the first-nearby settlement history contract by contract from official final CL settlements over the same January 2024--September 2026 inference window. The mapping contains 33 contract switches. Every return spanning a switch is excluded, and four isolated mapped settlement gaps are left missing rather than imputed; any one-day return that depends on a missing settlement is excluded as well.

The reconstructed series contains 632 usable daily returns, versus 676 for Yahoo. Its annualized realized volatility is 0.4183, compared with 0.4167 for Yahoo. Under the same Gaussian likelihood and prior, the posterior mean of $\sigma$ is 0.4182 with posterior standard deviation 0.0117, versus 0.4165 and 0.0112 under Yahoo. After normalizing Yahoo timestamps to trading dates, the 632 overlapping usable returns have correlation 0.9935 and mean absolute return difference 0.00051. Thus the two physical-volatility inputs are empirically very similar away from the deliberately excluded roll transitions.

The pricing result is equally stable. On exactly the same 2,381 strict-forward holdouts, replacing Yahoo with the reconstructed first-nearby history changes the historical posterior-integrated MAE/RMSE from 2.6187/3.4090 to 2.6468/3.4451. The expanding and previous-day option-informed errors remain 0.1088/0.1725 and 0.1075/0.1594 because their risk-neutral states are unchanged. On the 253-observation external-LO common-support sample, the reconstructed historical baseline has MAE 0.4362 and RMSE 0.5466, while the external previous-day LO surface remains at 0.1559 and 0.2599. The valuation-date cluster bootstrap still separates external LO from reconstructed historical volatility: the MAE difference is $-0.2804$ with a 95\% interval of $[-0.3329,-0.2243]$, and the RMSE difference is $-0.2867$ with interval $[-0.3399,-0.2490]$.

This check addresses the specific concern that the comparison between the long-window historical benchmark and option-informed inputs is an artifact of Yahoo's continuous-contract construction. It is not: the physical posterior barely moves and the pricing comparison is essentially unchanged.

\subsection{Numerical robustness: Monte Carlo, Sobol, and Curran}\label{sec:numerical_robustness}

The empirical PI--PM gaps are sufficiently small that numerical error is a first-order methodological concern. The high-precision benchmark therefore recomputes difficult contracts using three distinct engines. Pseudo-random Monte Carlo uses antithetic paths and an arithmetic-average control variate with known $Q$ expectation. Randomized Sobol uses independently scrambled low-discrepancy sequences without that control variate. Curran conditioning supplies a deterministic approximation under the same one-factor futures dynamics.

Across the difficult targets, the three engines agree on PI--PM to within a few micro-dollars. Direct seed-to-seed Monte Carlo uncertainty for the PI--PM gap is generally of order $10^{-6}$, compared with empirical gaps near $10^{-3}$; randomized-Sobol scramble uncertainty is of similarly smaller order. The conclusion that posterior integration adds little to the October point-price mean is therefore numerically identified rather than an artifact of insufficient simulation.

For representative October calls, the August 27, 2026 $K=98$ contract has a pseudo-Monte Carlo PI--PM gap of 0.0013683, a Curran gap of 0.0013690, and a randomized-Sobol gap of 0.0013717. The August 28 $K=98$ call gives 0.0013394, 0.0013402, and 0.0013393, respectively. Figure~\ref{fig:numerical_identification} shows all eight difficult empirical targets common to the three engines.

\begin{figure*}[t]
\centering
\includegraphics[width=\textwidth]{figures/publication/fig04_numerical_identification.pdf}
\caption{Numerical identification of the posterior-integration effect on eight difficult WTI APO targets. Panel A compares the PI--PM price difference from high-precision pseudo-random Monte Carlo, Curran conditioning, and independently randomized Sobol pricing; error bars show 95\% Monte Carlo or scramble-based intervals where applicable. Panel B reports deviations of the stochastic estimates from the deterministic Curran gap in units of $10^{-6}$ price units. Agreement at this scale shows that the observed sub-cent PI--PM differences are resolved well above numerical noise, even though the three engines need not produce identical absolute option prices.}
\label{fig:numerical_identification}
\end{figure*}

Curran prices can differ from Monte Carlo point prices by small amounts that are larger than the PI--PM correction itself. This is another useful distinction: pricing-engine approximation is a separate model/numerical margin from posterior integration. Importantly, Curran, Monte Carlo, and Sobol still produce almost the same \emph{difference} between PI and PM, so the integration conclusion does not depend on the chosen engine.

\subsection{Simulation-based calibration}

The Gaussian posterior implementation is validated with 200,000 prior-predictive datasets. After constructing the posterior CDF consistently with the trapezoidal quadrature density, posterior ranks are close to Uniform$(0,1)$ at every sample size.

For $n=21,63,252,$ and $1260$, rank means are 0.5000, 0.4993, 0.4978, and 0.5027, respectively; the target is 0.5. Rank variances are all close to $1/12$. Kolmogorov--Smirnov $p$-values are 0.861, 0.348, 0.210, and 0.106, so the experiment does not detect systematic rank non-uniformity. Empirical 95\% coverage is 0.9500, 0.9505, 0.9483, and 0.9505, with similarly accurate 50\% and 80\% coverage. Mean posterior bias in $\sigma$ is near zero. The calibration experiment therefore supports the statistical implementation used by the mechanism map and baseline posterior calculations.

\subsection{Liquidity and valuation-date dependence}

Option rows from the same date are not independent because they share the futures curve, historical posterior, and market environment. The robustness analysis therefore resamples valuation-date clusters rather than individual option contracts. In the extended forward experiment, the pooled sample contains 2,381 option-date observations over 15 distinct valuation-date clusters and seven expiries.

For the expanding prior-date smile, the pooled MAE difference relative to the historical-volatility posterior-integrated baseline is $-2.5099$, with a 95\% valuation-date cluster-bootstrap interval of $[-2.8297,-2.3000]$; the RMSE difference is $-3.2365$, with interval $[-3.5345,-3.0342]$. The previous-day smile produces corresponding differences of $-2.5112$ and $-3.2496$. Every resampled date-cluster draw in the reported bootstrap has lower error for these comparisons. This is a resampling result conditional on the observed dates, not a probability statement about all future market histories. The same direction survives open-interest thresholds of 10, 100, and 500.

The transaction-active pool is now materially larger than in the original September--October experiment. Across 62 positive-volume holdouts over 13 valuation-date clusters, the expanding smile lowers MAE by 1.5337 with a 95\% interval of $[-1.9039,-1.0401]$ and lowers RMSE by 2.2034 with interval $[-2.7027,-1.4213]$. The previous-day specification gives nearly identical conclusions. Thus the pooled positive-volume evidence now supports the same direction as the broader cross section. This does not eliminate the liquidity caveat: positive-volume observations remain sparse within individual long-dated expiries, especially in 2028, so maturity-specific inference is better supported by the denser open-interest-filtered panels than by transaction counts alone.

The expiry-level bootstrap reinforces that the pooled result is not generated by one maturity. For the expanding smile, the MAE improvement interval excludes zero for every expiry from September 2026 through September 2028. The absolute improvement grows with maturity because the historical volatility near 0.416 increasingly exceeds the lower effective APO volatility embedded in the 2027--2028 cross sections. This maturity pattern is economically informative but should not be read as a structural estimate of a long-run risk-neutral diffusion coefficient: stale marks, term-structure effects, smile dynamics, and model misspecification can all be absorbed by the one-factor implied-volatility state.

Cross-sectional calibration still reveals an important limitation. Leave-one-contract-out APO calibration strongly improves the broad September--October sample, while some narrow transaction-active and cross-option-type subsets are less stable. Those reversals are consistent with a non-flat implied-volatility surface and with calls and puts occupying different moneyness regions. They remain evidence against interpreting one common $\sigma_Q$ as a complete structural model of the APO cross-section.

\subsection{Interpretation of option-informed risk-neutral volatility}

The APO-implied exercise is intentionally endogenous to the option market being studied, so by itself it cannot show that the forward improvement survives outside the target family. The independent LO experiment supplies that missing check. It constructs the risk-neutral state from standard monthly WTI options and matching CL futures settlements, uses only dates strictly earlier than each APO target date, and then transfers the resulting surface to the APO fixing horizon. This design mirrors the market-calibration logic of \cite{shirayatakahashi2011} at a deliberately simpler model level: vanilla WTI options supply risk-neutral information before the average option is valued.

The external result strengthens the interpretation of the forward-smile evidence. On the 253 no-clipping common-support contract-dates, the previous-day LO surface has RMSE 0.2599 versus 0.5562 for historical PI and 0.2686 for the previous-day APO smile. The expanding LO and APO surfaces are similarly close. Valuation-date cluster bootstraps separate either option-informed specification from the long-window historical baseline, while most LO-versus-APO differences are small. The empirical conclusion is therefore not that one option family supplies a universally superior volatility estimate, nor that the experiment isolates a pure structural $\mathbb P$--$\mathbb Q$ wedge. It shows that the pricing advantage of recent option-informed volatility survives a cross-instrument construction.

Strike dependence remains important. The initial LO pilot covered only strikes 85--94.5 and placed just 60 of 280 matched contract-dates inside genuine prior moneyness support. Expanding the observed strike range to 70--120 raises common support to 253 observations, or 90.4\% of the matched sample, and reduces clipped external-surface rows to 50 of 620. This progression shows why a scalar or narrowly supported volatility comparison is insufficient. The result should still not be reduced to a claim that risk-neutral volatility is simply higher than physical volatility: the APO maturity structure changes sign relative to the historical level, and both the LO and APO cross sections contain smile information. The relevant object is an effective volatility surface under the maintained pricing map, not a unique structural diffusion coefficient.

\subsection{Remaining model and data limitations}

Several limitations remain relevant to the scope of the conclusions. The main historical tables retain Yahoo \texttt{CL=F} as the baseline return source for comparability, although the contract-reconstructed first-nearby robustness above shows that the empirical hierarchy is insensitive to that choice. The four isolated missing settlements in the reconstructed series are not imputed and all affected one-day returns are excluded. In the Barchart histories used here, \texttt{Latest} is the CME settlement field. Discounting uses an interpolated Treasury par-yield approximation rather than a bootstrapped OIS curve.

The pricing dynamics remain one-factor and constant-volatility. A correlated multi-factor futures curve, stochastic volatility, jumps, or mean reversion could improve absolute fit and would introduce additional uncertain parameters. Such models are natural extensions, but they should be evaluated separately from the narrower question answered here: conditional on a specified pricing map, how much does posterior integration add beyond a point estimate?

Finally, the forward-validation panel spans seven expiries from September 2026 through September 2028 but only 15 distinct target dates, while the independent vanilla experiment is concentrated on the October-2026 fixing structure and 10 matched APO valuation dates. Date-cluster resampling preserves within-date dependence but does not by itself remove serial dependence across adjacent dates. Liquidity also remains uneven: positive daily volume is sparse in the longest-dated APO expiries even when open interest is nonzero. The long-maturity evidence should therefore be interpreted as validation against available end-of-day market marks rather than as a dense transaction-price study. In the LO inversion panel, 5,214 of 5,220 settlements are flagged actual and six theoretical; this distinction is retained as a data-quality diagnostic rather than hidden. Broader cross-expiry vanilla validation would further strengthen external validity.

\section{Conclusion}\label{sec:conclusion}

This paper isolates a narrow Bayesian pricing question: how much does integrating uncertainty in historically estimated volatility change a WTI Average Price Option value relative to plugging in the posterior mean? The answer is governed by posterior dispersion and pricing curvature. Synthetic experiments confirm this mechanism and show that posterior integration can become economically relevant when information is weak and the pricing map is sufficiently curved.

The observed WTI sample lies in a different region. Posterior-integrated and posterior-mean point values are very close, yet the posterior distribution of theoretical prices remains materially wider. The distinction matters: a small \(C_{PI}-C_{PM}\) is a statement about nonlinear averaging of the posterior, not about the absence of parameter-induced price uncertainty.

The larger empirical margin is the volatility information supplied to the pricing map. Across the tested forward benchmarks, prior option-implied volatility states predict APO settlements much better than the historical-volatility inputs. That finding survives recent-window, EWMA, and horizon-matched GARCH comparisons; it does not require flexible smile fitting; and it remains visible when the volatility source is moved to an independent vanilla-WTI option market. These results support an option-informed versus tested-historical comparison, not a pure structural decomposition between physical and risk-neutral volatility.

The contribution is therefore diagnostic rather than a new option-pricing theory. Posterior integration matters when dispersion meets curvature, but in this sample improving the volatility information is much more consequential for point pricing than more exact integration of the historical-volatility posterior. The evidence is limited by the small number of distinct market dates, uneven liquidity, and the maintained one-factor pricing model. Richer commodity dynamics can be studied from this benchmark without conflating model enrichment with the narrower effect of Bayesian posterior integration.


\bmsection*{Data availability}
The APO and CL histories used in the empirical application were obtained through a commercial Barchart subscription. The independent vanilla-WTI validation uses CME/NYMEX LO option and CL futures statistics accessed through Databento. Redistribution of raw proprietary market-data files is not assumed to be permitted. The public repository therefore prioritizes code, query metadata, schemas, provenance records, diagnostics, hashes, and derived results consistent with the applicable data licenses.

\bmsection*{Code availability}
The research code is organized as an installable Python package, \texttt{bayesian\_asian\_options}, with documented APIs, tests, deterministic seeds, and reproducible experiment entry points. A versioned code snapshot will accompany the final article.

\bmsection*{Conflict of interest}
The author declares no conflict of interest.

\bmsection*{Acknowledgments}
To be completed before submission.

\bibliography{references}

\end{document}
